% ── chemistry-handout-preamble.tex ── 化学竞赛讲义 LaTeX 导言区（v7 · 精排版）
\documentclass[a4paper,11pt]{ctexart}

% ── 字体优化 ──
\setCJKmainfont[AutoFakeBold=true,BoldFont=SimHei]{SimSun}
\setCJKsansfont{SimHei}

% 化学式宏（mhchem）
\usepackage[version=4]{mhchem}

% ── 数学 ──
\usepackage{amsmath,amssymb}

% ── 图片 ──
\usepackage{graphicx}

% ── 表格（跨页+自动列宽）──
\usepackage{booktabs}
\usepackage{xltabular}
\usepackage{array}
\newcolumntype{L}{>{\raggedright\arraybackslash}X}
\newcolumntype{C}{>{\centering\arraybackslash}X}
\newcolumntype{R}{>{\raggedleft\arraybackslash}X}

% ── 浮动体（[H] 强制定位）──
\usepackage{float}

% ── 页面 ──
\usepackage[includeheadfoot,margin=2.2cm,headheight=14pt]{geometry}
\setlength{\textwidth}{16.6cm}

% ── 行间距 ──
\usepackage{setspace}
\setstretch{1.15}
\setlength{\parskip}{0.3em}
\setlength{\parindent}{2em}

% ── 防止表格跨页 ──
\usepackage{longtable}
\setlength{\LTpre}{0pt}
\setlength{\LTpost}{0pt}

% ── 防止图片跨页（尽量）──
\renewcommand{\textfraction}{0.1}
\renewcommand{\topfraction}{0.9}
\renewcommand{\bottomfraction}{0.9}
\setcounter{totalnumber}{10}
\setcounter{topnumber}{10}
\setcounter{bottomnumber}{10}

% ── 超链接(无红框) ──
\usepackage[colorlinks=true,linkcolor=black,citecolor=black,urlcolor=black]{hyperref}

% ── 颜色 ──
\usepackage[dvipsnames,svgnames]{xcolor}
\definecolor{chapterblue}{RGB}{23,50,77}
\definecolor{insightbg}{RGB}{235,240,250}
\definecolor{warningbg}{RGB}{255,245,235}
\definecolor{tipbg}{RGB}{235,250,235}

% ── 页眉页脚 ──
\usepackage{fancyhdr}
\pagestyle{fancy}
\fancyhf{}
\fancyhead[L]{\small\color{chapterblue}\leftmark}
\fancyhead[R]{\small\color{chapterblue}\thepage}
\renewcommand{\headrule}{{\color{chapterblue}\hrule height 0.8pt}}

% ── 章节编号与样式 ──
\setcounter{secnumdepth}{1}
\ctexset{
  section = {
    format = \Large\bfseries\heiti\color{chapterblue},
    aftername = \hspace{0.5em},
    beforeskip = 2ex,
    afterskip = 1ex,
  },
  subsection = {
    format = \large\bfseries\heiti,
    beforeskip = 1.5ex,
    afterskip = 0.5ex,
  },
  subsubsection = {
    format = \normalsize\bfseries\heiti,
    beforeskip = 1ex,
    afterskip = 0.3ex,
  },
  paragraph = {
    format = \small\bfseries,
    beforeskip = 0.5ex,
    afterskip = 0.2ex,
  },
}

% ── blockquote（引用/提示框）缩减 ──
\usepackage{etoolbox}
\AtBeginEnvironment{quote}{\small\setstretch{1.05}\leftskip=1em\rightskip=1em}

% ── 教学提示框（浅底色版）──
\newenvironment{insightbox}{%
  \par\smallskip
  \noindent{\small\bfseries\color{RoyalBlue}🧠 认知冲突}\par
  \leavevmode\colorbox{insightbg}{\parbox{\dimexpr\textwidth-2\fboxsep\relax}{%
  \ignorespaces}}%
}{\par\smallskip}

\newenvironment{warningbox}{%
  \par\smallskip
  \noindent{\small\bfseries\color{BrickRed}⚠️ 易错点}\par
  \leavevmode\colorbox{warningbg}{\parbox{\dimexpr\textwidth-2\fboxsep\relax}{%
  \ignorespaces}}%
}{\par\smallskip}

\newenvironment{tipbox}{%
  \par\smallskip
  \noindent{\small\bfseries\color{ForestGreen}💡 理解提示}\par
  \leavevmode\colorbox{tipbg}{\parbox{\dimexpr\textwidth-2\fboxsep\relax}{%
  \ignorespaces}}%
}{\par\smallskip}

% ── 列表（答案区加间距）──
\usepackage{enumitem}
\setlist{nosep}
% 专门用于参考答案的列表（每项之间空一行）
\newlist{answerlist}{enumerate}{1}
\setlist[answerlist]{label=\arabic*., leftmargin=2em, itemsep=0.8em}

% ── 防止内容溢出 ──
\usepackage{microtype}
\setlength{\emergencystretch}{2em}
\setlength{\hfuzz}{2pt}
\setlength{\overfullrule}{0pt}
\sloppy

% ── pandoc 兼容 ──
\providecommand{\tightlist}{\setlength{\itemsep}{0pt}\setlength{\parskip}{0pt}}
\newenvironment{Shaded}{}{}
\newenvironment{Highlighting}{}{}
\newcommand{\NormalTok}[1]{#1}
\newcommand{\CommentTok}[1]{#1}
\newcommand{\KeywordTok}[1]{#1}
\newcommand{\StringTok}[1]{#1}
\newcommand{\FunctionTok}[1]{#1}
\newcommand{\DataTypeTok}[1]{#1}
\newcommand{\NumberTok}[1]{#1}
\newcommand{\ControlFlowTok}[1]{#1}
\newcommand{\OperatorTok}[1]{#1}
\newcommand{\BaseNTok}[1]{#1}
\newcounter{none}

% ── 图片尺寸限制 ──
\AtBeginDocument{\setkeys{Gin}{width=\textwidth,keepaspectratio}}

% ── 防止图片跨节 ──
\usepackage{placeins}

% ── 自定义命令 ──
% 练习题标题（去掉括号）
% 例题标题（去掉中括号）
% 这些需要在convert.py预处理中做正则替换

% ── Unicode 符号渲染命令（preprint 中调用，\ensuremath 避免 $...$ 切断 CJK 字间距）──
\newcommand{\箭}{\ensuremath{\rightarrow}}
\newcommand{\左箭}{\ensuremath{\leftarrow}}
\newcommand{\上箭}{\ensuremath{\uparrow}}
\newcommand{\下箭}{\ensuremath{\downarrow}}
\newcommand{\衡}{\ensuremath{\rightleftharpoons}}
\newcommand{\双箭}{\ensuremath{\leftrightarrow}}
\newcommand{\约等于}{\ensuremath{\approx}}
\newcommand{\小于等于}{\ensuremath{\leq}}
\newcommand{\大于等于}{\ensuremath{\geq}}
\newcommand{\不等于}{\ensuremath{\neq}}
\newcommand{\西格马}{\ensuremath{\sigma}}
\newcommand{\派}{\ensuremath{\pi}}

\graphicspath{{C:/Obsidion/妙妙屋/11-模板/scripts/.chem_media/}}
\title{C:\Obsidion\妙妙屋\04-课件\学生讲义\热力学初步}
\date{}
\setlength{\tabcolsep}{3.5pt}
\small
\begin{document}
\maketitle
\tableofcontents
\newpage
\section{热力学初步}\label{ux70edux529bux5b66ux521dux6b65}

\subsection{学习目标}\label{ux5b66ux4e60ux76eeux6807}

\begin{itemize}
\tightlist
\item
  能区分体系与环境、状态函数与途径函数，正确判断恒压过程中 Q、W、ΔU、ΔH 的符号
\item
  能用盖斯定律、标准生成焓、键焓三种方法计算标准反应焓变 \(\Delta_r H^\theta\)
\item
  能区分 ΔU 与 ΔH 的物理意义，掌握 ΔH = ΔU + ΔnRT 的推导与应用
\item
  能定性判断化学反应熵变 ΔS 的符号，知道标准熵的概念与基本计算
\item
  能用 ΔG = ΔH − TΔS 判断反应自发性，识别四类反应并估算转变温度
\item
  能区分 ΔG 与 \(\Delta G^\theta\) 的物理意义，会用 \(\Delta G^\theta = -RT \ln K^\theta\) 建立热力学与化学平衡的桥梁
\item
  知道 \(\Delta G^\theta\) 与标准电动势 \(E^\theta\) 的关系(\(\Delta G^\theta = -nFE^\theta\))，为电化学学习做铺垫
\end{itemize}

\subsection{一、为什么需要热力学}\label{ux4e00ux4e3aux4ec0ux4e48ux9700ux8981ux70edux529bux5b66}

化学热力学的重要性不仅在于可以应用它的基本原理解释许多化学现象，而且还能依据这些原理去判断反应进行的方向，预测反应发生的可能性。

在日常化学工作中我们总会遇到这样或那样的实际问题，其中有许多要借助于热力学方法才能得到解决。例如：

\begin{itemize}
\tightlist
\item
  高炉炼铁过程中的主要反应 \(\mathrm{Fe_2O_3 + 3CO \to 2Fe + 3CO_2}\)，那能不能用类似方法炼铝？---热力学告诉我们：不能，Al\textsubscript{2}O\textsubscript{3} 的 \(\Delta_f H^\theta\) 太负
\item
  NO 和 CO 都是汽车尾气中的有毒成分，它们能否相互反应生成无毒的 N\textsubscript{2} 和 CO\textsubscript{2}？---热力学告诉我们：能，\(\Delta G^\theta\) \textless{} 0
\item
  天然金刚石非常珍贵，但其同素异形体石墨却极其普通和价廉，能否找到一种条件使石墨转化为金刚石？---热力学说可以，但需要极端条件
\end{itemize}

化学热力学就是解决这类''反应能不能发生、能放/吸多少热''问题的工具。它回答的是\textbf{可能性}问题，而不是\textbf{速率}问题---后者属于化学动力学的范畴。

\subsection{二、量热与反应热}\label{ux4e8cux91cfux70edux4e0eux53cdux5e94ux70ed}

\subsubsection{2.1 保温杯式量热计}\label{ux4fddux6e29ux676fux5f0fux91cfux70edux8ba1}

化学反应的热效应可以通过量热计直接测量。最简单的装置是\textbf{保温杯式量热计}---由保温和测温两部分组成。保温的目的是保证反应过程中不与外界发生热的交换。镀银的玻璃瓶胆或聚苯乙烯泡沫杯都具有这种绝热保温的性能。为了准确测量温度变化，需要使用比较精确的温度计，至少刻度是 0.1°C 的温度计，借助放大镜可以读出 0.01°C 的温度变化。

利用这种简单的仪器装置，可以测量中和热、溶解热以及其他溶液反应的热效应。反应所放热量 Q 应该等于量热计和反应后溶液升温所需的热量：

\[
Q = c V \rho \Delta t + C \Delta t = (c V \rho + C) \times \Delta t
\]

其中 \(\Delta t\) 是溶液温度升高值(\(\Delta t = t_{\text{终}} - t_{\text{始}}\))；\(c\) 为溶液的比热容；\(V\) 为反应后溶液的总体积；\(\rho\) 为溶液的密度；\(C\) 叫做量热计常数，它代表量热计各部件(杯体、搅拌器、加热器、温度计等)热容量之总和，即量热计每升高 1°C 所需的热量。

确定常数 C 最简便的方法是将一份冷水 \((t_1, V_1)\) 置于量热计中，随即再将另一份热水 \((t_2, V_2)\) 倒入，搅拌并记录最终温度 \(t\)。热水所放热量必定等于冷水所吸热量和量热计所吸热量之和，\(Q_{\text{放}} = Q_{\text{吸}}\)。

\textbf{例1(量热法测中和热)}：利用保温杯式量热计测定摩尔中和热。

\begin{enumerate}
\def\labelenumi{(\arabic{enumi})}
\item
  量热计常数测定：量取温度分别为 22.10°C 的冷水和 34.55°C 的热水各 100 cm\textsuperscript{3}，倒入保温杯中，测得混合后体系温度为 28.00°C。
\item
  分别量取 1.01 mol·dm\textsuperscript{-3} 的 HCl 溶液和 1.03 mol·dm\textsuperscript{-3} 的 NaOH 溶液各 100 cm\textsuperscript{3}，倒入保温杯中，两溶液起始温度均为 t\textsubscript{1} = 22.30°C。反应后体系温度升至 t\textsubscript{2} = 29.00°C。已知水的密度为 1.00 g·cm\textsuperscript{-3}，比热容 c = 4.03 J·g\textsuperscript{-1}·°C\textsuperscript{-1}，实验所得 NaCl 溶液密度 ρ = 1.02 g·cm\textsuperscript{-3}。
\end{enumerate}

解：先由冷水和热水温度变化求量热计常数 C：

\[(34.55 - 28.00) \times 100 \times 1.00 \times 4.18 = (28.00 - 22.10) \times (100 \times 1.00 \times 4.18 + C)\]

求得 C = 46.1 J·°C\textsuperscript{-1}。

代入量热公式：\(Q = (200 \times 1.02 \times 4.03 + 46.1) \times (29.00 - 22.30) = 5.82\) kJ

摩尔中和热 = 5.82 / (1.01 × 0.100) = \textbf{57.6 kJ·mol\textsuperscript{-1}}

查阅化学手册可知，1 mol H\textsuperscript{+} 和 1 mol OH\textsuperscript{-} 反应放热 57.3 kJ，实验结果与此数据相符。

\subsubsection{2.2 弹式量热计}\label{ux5f39ux5f0fux91cfux70edux8ba1}

弹式量热计(氧弹量热计)适用于测定燃烧热。化学反应在一个厚壁钢制容器内进行，该容器形状像小炸弹。在实验前必须向钢弹中通入一定量的高压氧气。弹盖由细密螺纹旋紧，整个氧弹位于有绝热外套的水浴之中。钢弹是密闭容器，反应过程中总体积不变，测定的是\textbf{恒容反应热} \(Q_V\)：

\[
Q_V = Q_{\text{水}} + Q_{\text{弹}} = 4.18 \times m \times \Delta t + C \times \Delta t
\]

常用标准物是苯甲酸(C\textsubscript{6}H\textsubscript{5}COOH)，其摩尔燃烧热为 3.23 × 10\textsuperscript{3} kJ·mol\textsuperscript{-1}，用来标定量热计常数 C。

\textbf{理解要点}：保温杯式量热计测的是恒压反应热 \(Q_p = \Delta H\)，弹式量热计测的是恒容反应热 \(Q_V = \Delta U\)。两者的关系为 \(\Delta H = \Delta U + \Delta nRT\)。

\subsection{三、内能与焓}\label{ux4e09ux5185ux80fdux4e0eux7113}

\subsubsection{3.1 热力学能(内能)U}\label{ux70edux529bux5b66ux80fdux5185ux80fdu}

内能是物质的一种属性，指体系内物质所含分子及原子的动能、势能、核能、电子能等能量的总和。内能不包括体系宏观运动的动能和体系在外力场中的势能。

内能是一种\textbf{状态函数}---物质处于一定的状态，就具有一定的内能，状态发生改变，内能也随之改变。物质经过一系列变化之后，如果又回到原来的状态，内能也回复原值。但一定状态下，物质内能的绝对值\textbf{无法直接测定}，内能的变化 \(\Delta U\) 可以通过测定过程的热和功而得到。

\subsubsection{3.2 热力学第一定律}\label{ux70edux529bux5b66ux7b2cux4e00ux5b9aux5f8b}

能量守恒定律：能量可以从一种形式转化为另一种形式，但总值不变。

\textbf{数学表达式}：

\[
\Delta U = Q + W
\]

其中 \(\Delta U = U_2 - U_1\) 是内能变化；\(Q\) 是体系吸收或放出的热；\(W\) 是体系对环境或环境对体系所做的功。

{\def\LTcaptype{none} % do not increment counter
\begin{longtable}[]{@{}cll@{}}
\toprule\noalign{}
物理量 & \textgreater{} 0 & \textless{} 0 \\
\midrule\noalign{}
\endhead
\bottomrule\noalign{}
\endlastfoot
ΔU & 体系内能增加 & 体系内能减少 \\
Q & 体系吸收热量 & 体系放出热量 \\
W & 环境对体系做功 & 体系对环境做功 \\
\end{longtable}
}

\textbf{易错点(物理 vs 化学符号约定)}：物理教材中 \(\Delta U = Q - W\)(W 定义为体系对外做功)，化学教材中 \(\Delta U = Q + W\)(W 定义为环境对体系做功)。两者不矛盾---只是 W 的取号约定不同。化学约定中，环境对体系做功 W \textgreater{} 0，体系对环境做功 W \textless{} 0。

\subsubsection{3.3 焓 H}\label{ux7113-h}

\textbf{定义}：

\[
H \equiv U + pV
\]

焓是状态函数。对于\textbf{恒压且只做体积功}的过程，由第一定律：

\[
\Delta U = Q_p + W = Q_p - p\Delta V
\]

展开得 \((U_2 + pV_2) - (U_1 + pV_1) = Q_p\)

即：

\[
\boxed{\Delta H = Q_p}
\]

即恒压过程中的热效应等于体系焓的变化。当化学反应在恒压下进行，所测定的反应热为恒压反应热 \(Q_p\)。

\textbf{理解要点}：焓 H 是人们在处理体系状态变化时引入的一个状态函数。一个化学反应是吸热还是放热，在等压条件下是由生成物和反应物的焓值之差决定的：\(\sum H(\text{生成物}) - \sum H(\text{反应物}) = \Delta H\)。

\subsubsection{3.4 ΔU 与 ΔH 的关系}\label{ux3b4u-ux4e0e-ux3b4h-ux7684ux5173ux7cfb}

将恒压条件代入第一定律：

\[
\Delta U = \Delta H - p\Delta V
\]

对气相反应，由理想气体方程 \(p\Delta V = \Delta nRT\)(\(\Delta n\) 为气态生成物与气态反应物的物质的量之差)：

\[
\boxed{\Delta H = \Delta U + \Delta nRT}
\]

对固/液相反应，\(\Delta V \approx 0\)，故 \(\Delta H \approx \Delta U\)。

\textbf{例2(ΔU 与 ΔH 的换算)}：1.00 g 联氨(N\textsubscript{2}H\textsubscript{4}，M = 32.0 g/mol)在氧气中完全燃烧(等容)时放热 20.7 kJ。求 1 mol N\textsubscript{2}H\textsubscript{4} 燃烧时的内能变化和等压反应热。

解：反应方程式 \(\mathrm{N_2H_4(g) + O_2(g) \to N_2(g) + 2H_2O(l)}\)

\[Q_V = -20.7 \times 32.0 = -662\ \mathrm{kJ/mol}, \quad \Delta U = Q_V = -662\ \mathrm{kJ/mol}\]

反应前后气态物质变化：\(\Delta n = 1 - 2 = -1\)

\[\Delta H = \Delta U + \Delta nRT = -662 + (-1) \times 8.314 \times 10^{-3} \times 298 = -665\ \mathrm{kJ/mol}\]

答案：\(\Delta U = -662\) kJ/mol，\(\Delta H = -665\) kJ/mol。

\subsection{四、热化学方程式与盖斯定律}\label{ux56dbux70edux5316ux5b66ux65b9ux7a0bux5f0fux4e0eux76d6ux65afux5b9aux5f8b}

\subsubsection{4.1 热化学方程式}\label{ux70edux5316ux5b66ux65b9ux7a0bux5f0f}

表示化学反应与热效应关系的方程式称为\textbf{热化学方程式}：

\[
\mathrm{H_2(g) + \frac{1}{2}O_2(g) = H_2O(l)}, \quad \Delta_r H^\theta = -285.8\ \mathrm{kJ \cdot mol^{-1}}
\]

\textbf{书写规则}：

{\def\LTcaptype{none} % do not increment counter
\begin{longtable}[]{@{}
  >{\raggedright\arraybackslash}p{(\linewidth - 4\tabcolsep) * \real{0.3333}}
  >{\raggedright\arraybackslash}p{(\linewidth - 4\tabcolsep) * \real{0.3333}}
  >{\raggedright\arraybackslash}p{(\linewidth - 4\tabcolsep) * \real{0.3333}}@{}}
\toprule\noalign{}
\begin{minipage}[b]{\linewidth}\raggedright
规则
\end{minipage} & \begin{minipage}[b]{\linewidth}\raggedright
说明
\end{minipage} & \begin{minipage}[b]{\linewidth}\raggedright
示例
\end{minipage} \\
\midrule\noalign{}
\endhead
\bottomrule\noalign{}
\endlastfoot
1. 标注物态 & 反应物和生成物须注明 (g)/(l)/(s)/(aq) 或晶型 & C(石墨) 与 C(金刚石) 不同 \\
2. ΔH 符号 & 放热 ΔH \textless{} 0，吸热 ΔH \textgreater{} 0 & 放热写''−``，吸热写''+'' \\
3. 系数为物质的量 & 可为分数，不代表分子个数 & 可写 \(\frac{1}{2}\mathrm{O_2}\) \\
4. ΔH 与系数成正比 & 系数加倍，ΔH 数值也加倍；反向则符号相反 & --- \\
\end{longtable}
}

\subsubsection{4.2 盖斯定律}\label{ux76d6ux65afux5b9aux5f8b}

在恒 T、p(或恒 T、V)条件下，化学反应的总热效应\textbf{只与始态和终态有关，与具体途径无关}。

\textbf{本质}：焓 H 是状态函数 \箭 ΔH 只取决于始态和终态。

\textbf{应用方法}：将目标反应拆分为若干已知 ΔH 的子反应，通过代数组合求算。

\textbf{例3(盖斯定律)}：已知 298 K 时：
- (1) C(s) + O\textsubscript{2}(g) = CO\textsubscript{2}(g)，\(\Delta_r H_1^\theta\) = −393.5 kJ/mol
- (2) H\textsubscript{2}(g) + ½O\textsubscript{2}(g) = H\textsubscript{2}O(l)，\(\Delta_r H_2^\theta\) = −285.8 kJ/mol
- (3) CH\textsubscript{4}(g) + 2O\textsubscript{2}(g) = CO\textsubscript{2}(g) + 2H\textsubscript{2}O(l)，\(\Delta_r H_3^\theta\) = −890.3 kJ/mol

求反应 C(s) + 2H\textsubscript{2}(g) = CH\textsubscript{4}(g) 的 \(\Delta_r H^\theta\)。

解：目标反应 = (1) + 2×(2) − (3)

\[\Delta_r H^\theta = (-393.5) + 2 \times (-285.8) - (-890.3) = -74.8\ \mathrm{kJ/mol}\]

答案：\(\Delta_r H^\theta = -74.8\) kJ/mol(这正是 CH\textsubscript{4} 的标准生成焓)

\subsubsection{4.3 三种焓变计算方法}\label{ux4e09ux79cdux7113ux53d8ux8ba1ux7b97ux65b9ux6cd5}

\textbf{(1) 标准生成焓 \(\Delta_f H^\theta\) 法(最常用)}

在标准状态和指定温度(通常 298.15 K)下，由\textbf{稳定单质}生成 1 mol 化合物时的焓变。稳定单质的 \(\Delta_f H^\theta\) = 0。

\[
\boxed{\Delta_r H^\theta = \sum n_i \Delta_f H^\theta(\text{生成物}) - \sum n_j \Delta_f H^\theta(\text{反应物})}
\]

注意：C 的稳定单质是\textbf{石墨}而非金刚石，\(\Delta_f H^\theta(\text{C, 金刚石}) = +1.9\) kJ/mol \不等于 0。

\textbf{例4(生成焓法)}：计算反应 CH\textsubscript{4}(g) + Cl\textsubscript{2}(g) = CH\textsubscript{3}Cl(g) + HCl(g) 的 \(\Delta_r H^\theta\)。

已知：\(\Delta_f H^\theta(\text{CH}_4) = -74.8\)，\(\Delta_f H^\theta(\text{CH}_3\text{Cl}) = -81.9\)，\(\Delta_f H^\theta(\text{HCl}) = -92.3\)，\(\Delta_f H^\theta(\text{Cl}_2) = 0\)(均为 kJ/mol)

\[\Delta_r H^\theta = [(-81.9) + (-92.3)] - [(-74.8) + 0] = -99.4\ \mathrm{kJ/mol}\]

\textbf{(2) 标准燃烧焓 \(\Delta_c H^\theta\) 法}

1 mol 物质完全燃烧的焓变。完全燃烧规定：C\箭CO\textsubscript{2}(g)，H\箭H\textsubscript{2}O(l)，N\箭N\textsubscript{2}(g)，S\箭SO\textsubscript{2}(g)，Cl\箭HCl(aq)。

\[
\Delta_r H^\theta = \sum n_i \Delta_c H^\theta(\text{反应物}) - \sum n_j \Delta_c H^\theta(\text{生成物})
\]

公式方向与生成焓法\textbf{相反}！因为燃烧焓是''物质被烧掉''的焓变。

\textbf{(3) 键焓 \(\Delta_b H^\theta\) 法(估算用)}

1 mol 气态物质断开化学键成为气态原子时的焓变。\(\Delta_r H^\theta\) = Σ(反应物断键) − Σ(生成物成键)。

\textbf{键焓法为什么是估算}：键焓是\textbf{平均值}(从多种分子中统计得到)，忽略了分子环境差异。例如不同分子中 C−H 键的实际键能略有不同，但键焓法统一用 413 kJ/mol。

{\def\LTcaptype{none} % do not increment counter
\begin{longtable}[]{@{}llcl@{}}
\toprule\noalign{}
方法 & 适用条件 & 精度 & 典型场景 \\
\midrule\noalign{}
\endhead
\bottomrule\noalign{}
\endlastfoot
生成焓法 & 有热力学数据表 & 高 & 无机反应、热化学计算 \\
燃烧焓法 & 有燃烧焓数据 & 高 & 有机反应(燃烧焓数据丰富) \\
键焓法 & 缺热力学数据但有键能表 & 低 & 估算、无数据时的快速判断 \\
\end{longtable}
}

\subsubsection{4.4 Born-Haber 循环}\label{born-haber-ux5faaux73af}

盖斯定律在离子晶体中的经典应用。以 NaCl 的形成为例：

\[\text{Na(s)} + \frac{1}{2}\text{Cl}_2\text{(g)} \to \text{NaCl(s)}\]

可以拆解为五个步骤：
1. Na(s) \箭 Na(g) \textbf{升华热} \(\Delta_{\text{sub}}H\) = +107 kJ/mol
2. Na(g) \箭 Na\textsuperscript{+}(g) + e\textsuperscript{-} \textbf{第一电离能} \(I_1\) = +496 kJ/mol
3. ½Cl\textsubscript{2}(g) \箭 Cl(g) \textbf{解离能} \(\frac{1}{2}D\) = +122 kJ/mol
4. Cl(g) + e\textsuperscript{-} \箭 Cl\textsuperscript{-}(g) \textbf{电子亲和能} \(E_{\text{ea}}\) = −349 kJ/mol
5. Na\textsuperscript{+}(g) + Cl\textsuperscript{-}(g) \箭 NaCl(s) \textbf{晶格能} \(U\)

由盖斯定律：\(\Delta_f H^\theta = \Delta_{\text{sub}}H + I_1 + \frac{1}{2}D + E_{\text{ea}} + U\)

Born-Haber 循环的价值：晶格能无法直接实验测量，但通过 Born-Haber 循环可以从其他可测量的量间接求出。

\subsection{五、熵与热力学第二定律}\label{ux4e94ux71b5ux4e0eux70edux529bux5b66ux7b2cux4e8cux5b9aux5f8b}

\subsubsection{5.1 为什么需要熵}\label{ux4e3aux4ec0ux4e48ux9700ux8981ux71b5}

仅用 ΔH 无法判断反应方向。例如 N\textsubscript{2}(g) + O\textsubscript{2}(g) = 2NO(g) 的 ΔH = +180.5 kJ/mol(吸热)，但这个反应在高温下确实能发生(闪电固氮)。说明除焓之外，还有另一个因素在驱动反应---\textbf{熵}。

\subsubsection{5.2 熵 S 的定义}\label{ux71b5-s-ux7684ux5b9aux4e49}

熵是衡量体系\textbf{混乱度}(无序度)的物理量，是\textbf{状态函数}。从微观角度看，熵是体系\textbf{微观状态数} Ω 的量度：

\[
S = k_B \ln \Omega
\]

其中 \(k_B\) 为 Boltzmann 常数(1.381 × 10\textsuperscript{-23} J/K)。

\textbf{热力学第三定律}：在 0 K 时，任何\textbf{完美晶体}的熵为零。由此可以计算物质在各温度下的\textbf{标准摩尔熵} \(S^\theta\)(规定熵)。

\textbf{标准反应熵变}：

\[
\Delta_r S^\theta = \sum n_i S^\theta(\text{生成物}) - \sum n_j S^\theta(\text{反应物})
\]

\textbf{关键区别}：与生成焓不同，\textbf{稳定单质的标准熵不为零}！因为 0 K 以上就有热运动。

\subsubsection{5.3 定性判断 ΔS 符号(六条经验规则)}\label{ux5b9aux6027ux5224ux65ad-ux3b4s-ux7b26ux53f7ux516dux6761ux7ecfux9a8cux89c4ux5219}

{\def\LTcaptype{none} % do not increment counter
\begin{longtable}[]{@{}lcl@{}}
\toprule\noalign{}
变化 & ΔS 方向 & 原因 \\
\midrule\noalign{}
\endhead
\bottomrule\noalign{}
\endlastfoot
固 \箭 液 \箭 气 & 增大 (+) & 分子运动自由度增大 \\
气体分子数增加 & 增大 (+) & 微观状态数增多 \\
温度升高 & 增大 (+) & 分子运动加剧 \\
体积增大(气体) & 增大 (+) & 分子可占据空间增大 \\
溶质溶解(多数) & 增大 (+) & 溶质分子分散 \\
气体被吸收/溶解 & 减小 (−) & 气体分子自由度降低 \\
\end{longtable}
}

记忆口诀：``固体变液体变气体，分子数增多体积增大，熵就增大。'' 从有序到无序，熵增大。

\subsection{六、Gibbs 自由能与反应方向}\label{ux516dgibbs-ux81eaux7531ux80fdux4e0eux53cdux5e94ux65b9ux5411}

\subsubsection{6.1 Gibbs 自由能 G}\label{gibbs-ux81eaux7531ux80fd-g}

\textbf{定义}：

\[
G \equiv H - TS
\]

G 是状态函数。在\textbf{恒温恒压}条件下，反应自发性的判据为：

\[
\boxed{\Delta G = \Delta H - T\Delta S}
\]

{\def\LTcaptype{none} % do not increment counter
\begin{longtable}[]{@{}cl@{}}
\toprule\noalign{}
条件 & 含义 \\
\midrule\noalign{}
\endhead
\bottomrule\noalign{}
\endlastfoot
\(\Delta G < 0\) & 自发过程(正向可行) \\
\(\Delta G = 0\) & 平衡状态 \\
\(\Delta G > 0\) & 非自发过程(逆向可行) \\
\end{longtable}
}

\subsubsection{6.2 四类反应的自发性}\label{ux56dbux7c7bux53cdux5e94ux7684ux81eaux53d1ux6027}

{\def\LTcaptype{none} % do not increment counter
\begin{longtable}[]{@{}
  >{\centering\arraybackslash}p{(\linewidth - 12\tabcolsep) * \real{0.1471}}
  >{\centering\arraybackslash}p{(\linewidth - 12\tabcolsep) * \real{0.1471}}
  >{\centering\arraybackslash}p{(\linewidth - 12\tabcolsep) * \real{0.1471}}
  >{\centering\arraybackslash}p{(\linewidth - 12\tabcolsep) * \real{0.1471}}
  >{\centering\arraybackslash}p{(\linewidth - 12\tabcolsep) * \real{0.1471}}
  >{\centering\arraybackslash}p{(\linewidth - 12\tabcolsep) * \real{0.1471}}
  >{\raggedright\arraybackslash}p{(\linewidth - 12\tabcolsep) * \real{0.1176}}@{}}
\toprule\noalign{}
\begin{minipage}[b]{\linewidth}\centering
类型
\end{minipage} & \begin{minipage}[b]{\linewidth}\centering
ΔH
\end{minipage} & \begin{minipage}[b]{\linewidth}\centering
ΔS
\end{minipage} & \begin{minipage}[b]{\linewidth}\centering
低温
\end{minipage} & \begin{minipage}[b]{\linewidth}\centering
高温
\end{minipage} & \begin{minipage}[b]{\linewidth}\centering
转变温度
\end{minipage} & \begin{minipage}[b]{\linewidth}\raggedright
典型实例
\end{minipage} \\
\midrule\noalign{}
\endhead
\bottomrule\noalign{}
\endlastfoot
1. 焓降熵增 & − & + & 自发 & 自发 & 无 & 燃烧反应、中和反应 \\
2. 焓增熵减 & + & − & 非自发 & 非自发 & 无 & 2CO \箭 C + O\textsubscript{2}(极罕见) \\
3. 焓增熵增 & + & + & 非自发 & 自发 & \(T_{\text{转}} = \Delta H/\Delta S\) & CaCO\textsubscript{3} 分解 \\
4. 焓降熵减 & − & − & 自发 & 非自发 & \(T_{\text{转}} = \Delta H/\Delta S\) & NH\textsubscript{3}(g) + HCl(g) \箭 NH\textsubscript{4}Cl(s) \\
\end{longtable}
}

\textbf{理解要点}：ΔG \textless{} 0 不意味着反应一定快---热力学只讨论''能不能''，不讨论''有多快''。例如 H\textsubscript{2} 与 O\textsubscript{2} 在室温下 ΔG 极负，但没有催化剂时反应极慢。

\subsubsection{\texorpdfstring{6.3 \(\Delta G^\theta\) 与化学平衡常数}{6.3 \textbackslash Delta G\^{}\textbackslash theta 与化学平衡常数}}\label{delta-gtheta-ux4e0eux5316ux5b66ux5e73ux8861ux5e38ux6570}

\[
\boxed{\Delta_r G^\theta = -RT \ln K^\theta}
\]

其中 \(K^\theta\) 为标准平衡常数(无量纲)。这就是连接热力学与化学平衡的\textbf{桥梁公式}。

{\def\LTcaptype{none} % do not increment counter
\begin{longtable}[]{@{}
  >{\centering\arraybackslash}p{(\linewidth - 4\tabcolsep) * \real{0.3571}}
  >{\centering\arraybackslash}p{(\linewidth - 4\tabcolsep) * \real{0.3571}}
  >{\raggedright\arraybackslash}p{(\linewidth - 4\tabcolsep) * \real{0.2857}}@{}}
\toprule\noalign{}
\begin{minipage}[b]{\linewidth}\centering
\(\Delta G^\theta\) 大小
\end{minipage} & \begin{minipage}[b]{\linewidth}\centering
\(K^\theta\) 范围
\end{minipage} & \begin{minipage}[b]{\linewidth}\raggedright
含义
\end{minipage} \\
\midrule\noalign{}
\endhead
\bottomrule\noalign{}
\endlastfoot
\(\Delta G^\theta\) ≪ 0(如 \textless{} −40 kJ) & \(K^\theta\) ≫ 1 & 反应正向进行彻底 \\
\(\Delta G^\theta\) \约等于 0 & \(K^\theta\) \约等于 1 & 反应物与产物共存 \\
\(\Delta G^\theta\) ≫ 0(如 \textgreater{} +40 kJ) & \(K^\theta\) ≪ 1 & 反应几乎不发生 \\
\end{longtable}
}

\textbf{ΔG 与 \(\Delta G^\theta\) 的关系}(非标准状态下)：

\[
\Delta G = \Delta G^\theta + RT\ln Q
\]

当反应商 \(Q < K^\theta\) 时，\(\Delta G < 0\)，正向自发。

\textbf{核心辨析：\(\Delta G^\theta\) \不等于 ΔG}
- \(\Delta G^\theta\) 是标准状态下的值(所有气体 \(p^\theta\) = 100 kPa，溶液 \(c^\theta\) = 1 mol/L)，是\textbf{定值}
- ΔG 是任意状态下的值，\textbf{随浓度/分压变化}
- \(\Delta G^\theta\) \textgreater{} 0 不代表反应在任何条件下都不能发生---调整 Q 可能使 ΔG \textless{} 0

\subsubsection{\texorpdfstring{6.4 \(\Delta G^\theta\) 与标准电动势 \(E^\theta\)}{6.4 \textbackslash Delta G\^{}\textbackslash theta 与标准电动势 E\^{}\textbackslash theta}}\label{delta-gtheta-ux4e0eux6807ux51c6ux7535ux52a8ux52bf-etheta}

\[
\Delta_r G^\theta = -nFE^\theta_{\text{cell}}
\]

其中 n 为转移电子数，F = 96,485 C·mol\textsuperscript{-1}(Faraday 常数)。详见 电化学基础-超级充实版(自学完整)\textbar 电化学基础。

\textbf{单位陷阱预警}：R = 8.314 J/(mol·K)，而 \(\Delta G^\theta\) 通常以 kJ/mol 为单位。计算时必须先统一单位！单位错误是热力学计算失败的第一大原因。

\subsection{七、典型例题}\label{ux4e03ux5178ux578bux4f8bux9898}

\textbf{题干}：CaCO\textsubscript{3} 分解反应 \(\mathrm{CaCO_3(s) = CaO(s) + CO_2(g)}\)，已知 298 K 时：

{\def\LTcaptype{none} % do not increment counter
\begin{longtable}[]{@{}lcc@{}}
\toprule\noalign{}
物质 & \(\Delta_f H^\theta\) (kJ/mol) & \(S^\theta\) (J·mol\textsuperscript{-1}·K\textsuperscript{-1}) \\
\midrule\noalign{}
\endhead
\bottomrule\noalign{}
\endlastfoot
CaCO\textsubscript{3}(s) & −1206.9 & 92.9 \\
CaO(s) & −635.1 & 39.7 \\
CO\textsubscript{2}(g) & −393.5 & 213.8 \\
\end{longtable}
}

\begin{enumerate}
\def\labelenumi{(\arabic{enumi})}
\tightlist
\item
  计算 298 K 时的 \(\Delta_r H^\theta\)、\(\Delta_r S^\theta\)、\(\Delta_r G^\theta\)
\item
  估算标准压力下自发进行的最低温度
\end{enumerate}

\textbf{解析}：

\begin{enumerate}
\def\labelenumi{(\arabic{enumi})}
\tightlist
\item
  \(\Delta_r H^\theta = [(-635.1) + (-393.5)] - (-1206.9) = +178.3\) kJ/mol
\end{enumerate}

\(\Delta_r S^\theta = (39.7 + 213.8) - 92.9 = +160.6\) J/(mol·K) = +0.1606 kJ/(mol·K)

\(\Delta_r G^\theta = 178.3 - 298 \times 0.1606 = +130.4\) kJ/mol \textgreater{} 0

298 K 时 \(\Delta G^\theta\) \textgreater{} 0，反应非自发。

\begin{enumerate}
\def\labelenumi{(\arabic{enumi})}
\setcounter{enumi}{1}
\tightlist
\item
  令 \(\Delta G^\theta\) = 0：\(T_{\text{转}} = \frac{\Delta H^\theta}{\Delta S^\theta} = \frac{178.3}{0.1606} \approx 1110\) K \约等于 837°C
\end{enumerate}

答案：(1) \(\Delta_r H^\theta\) = +178.3 kJ/mol，\(\Delta_r S^\theta\) = +160.6 J/(mol·K)，\(\Delta_r G^\theta\) = +130.4 kJ/mol；(2) 约 1110 K

\textbf{题干}：298 K 时 \(\mathrm{N_2O_4(g) \rightleftharpoons 2NO_2(g)}\) 的 \(\Delta_r G^\theta = +4.73\) kJ/mol。求 \(K^\theta\)。

\textbf{解析}：

\[\ln K^\theta = \frac{-\Delta_r G^\theta}{RT} = \frac{-4.73 \times 10^3}{8.314 \times 298} = -1.91\]

\[K^\theta = e^{-1.91} = 0.148\]

答案：\(K^\theta\) = 0.148(\textless{} 1，说明 298 K 时 N\textsubscript{2}O\textsubscript{4} 更稳定)

\textbf{题干}：计算 CH\textsubscript{4}(g) + Cl\textsubscript{2}(g) = CH\textsubscript{3}Cl(g) + HCl(g) 的 \(\Delta_r H^\theta\)。

已知生成焓数据同例4；键焓：C−H = 413，Cl−Cl = 243，C−Cl = 339，H−Cl = 431 kJ/mol

\textbf{解析}：

\begin{enumerate}
\def\labelenumi{(\arabic{enumi})}
\item
  生成焓法：\(\Delta_r H^\theta = -99.4\) kJ/mol(精确值)
\item
  键焓法(断键 − 成键)：断键 = 413 + 243 = 656 kJ；成键 = 339 + 431 = 770 kJ
\end{enumerate}

\[\Delta_r H^\theta = 656 - 770 = -114\ \mathrm{kJ/mol}\]

\textbf{反思}：两者差 15 kJ/mol(约 15\%)，因为键焓是\textbf{平均值}，忽略了分子环境差异。

\subsection{八、知识网络与方法论}\label{ux516bux77e5ux8bc6ux7f51ux7edcux4e0eux65b9ux6cd5ux8bba}

\subsubsection{8.1 热力学函数关系总图}\label{ux70edux529bux5b66ux51fdux6570ux5173ux7cfbux603bux56fe}

\begin{verbatim}
                    U（内能）
                   ↙        ↘
            H = U + pV      W = −pΔV
              ↓                ↓
         ΔH = Q_p        ΔU = Q + W
              ↓
    ΔG = ΔH − TΔS
       ↙         ↘
  $\Delta G^\theta = -nFE^\theta$    $\Delta G^\theta = -RT \ln K^\theta$
\end{verbatim}

\subsubsection{8.2 热化学计算三步法}\label{ux70edux5316ux5b66ux8ba1ux7b97ux4e09ux6b65ux6cd5}

\begin{enumerate}
\def\labelenumi{\arabic{enumi}.}
\tightlist
\item
  \textbf{写方程式}：配平 + 标注物态
\item
  \textbf{选方法}：生成焓法(最准)/ 燃烧焓法(有机)/ 键焓法(估算)
\item
  \textbf{代数据}：\(\sum(\text{生成物}) - \sum(\text{反应物})\)---注意燃烧焓法方向相反！
\end{enumerate}

\subsubsection{8.3 反应自发性判断流程}\label{ux53cdux5e94ux81eaux53d1ux6027ux5224ux65adux6d41ux7a0b}

\begin{verbatim}
已知 ΔH 和 ΔS 的符号？
├─ ΔH < 0, ΔS > 0 → 所有温度自发
├─ ΔH > 0, ΔS < 0 → 所有温度非自发
├─ ΔH > 0, ΔS > 0 → 高温自发（T > ΔH/ΔS）
└─ ΔH < 0, ΔS < 0 → 低温自发（T < ΔH/ΔS）
\end{verbatim}

\subsection{九、本节总结速查}\label{ux4e5dux672cux8282ux603bux7ed3ux901fux67e5}

{\def\LTcaptype{none} % do not increment counter
\begin{longtable}[]{@{}
  >{\raggedright\arraybackslash}p{(\linewidth - 4\tabcolsep) * \real{0.3333}}
  >{\raggedright\arraybackslash}p{(\linewidth - 4\tabcolsep) * \real{0.3333}}
  >{\raggedright\arraybackslash}p{(\linewidth - 4\tabcolsep) * \real{0.3333}}@{}}
\toprule\noalign{}
\begin{minipage}[b]{\linewidth}\raggedright
核心概念
\end{minipage} & \begin{minipage}[b]{\linewidth}\raggedright
关键公式
\end{minipage} & \begin{minipage}[b]{\linewidth}\raggedright
注意点
\end{minipage} \\
\midrule\noalign{}
\endhead
\bottomrule\noalign{}
\endlastfoot
内能 U & \(\Delta U = Q + W\) & 状态函数；绝对值不可测 \\
焓 H & \(H = U + pV\)，\(\Delta H = Q_p\) & 恒压反应热 \\
ΔH 与 ΔU & \(\Delta H = \Delta U + \Delta nRT\) & 气相反应需考虑 Δn \\
盖斯定律 & ΔH 与途径无关 & 代数组合已知反应 \\
生成焓法 & \(\Delta_r H^\theta = \sum n_i \Delta_f H^\theta(\text{生}) - \sum\) & 稳定单质 \(\Delta_f H^\theta\) = 0 \\
燃烧焓法 & \(\Delta_r H^\theta = \sum(\text{反应物}) - \sum(\text{生成物})\) & 方向与生成焓法相反 \\
键焓法 & \(\Delta_r H^\theta = \sum(\text{断键}) - \sum(\text{成键})\) & 估算值 \\
Born-Haber 循环 & 盖斯定律在离子晶体中的应用 & 可间接求晶格能 \\
熵 S & 混乱度的量度；\(S = k_B \ln \Omega\) & 稳定单质 \(S^\theta\) \不等于 0 \\
Gibbs 自由能 & \(G = H - TS\) & 恒温恒压自发性判据 \\
自发性 & \(\Delta G = \Delta H - T\Delta S\) & ΔG \textless{} 0 自发 \\
\(\Delta G^\theta\)---\(K^\theta\) & \(\Delta G^\theta = -RT\ln K^\theta\) & \(K^\theta\) 无量纲；单位统一 \\
\(\Delta G^\theta\)---\(E^\theta\) & \(\Delta G^\theta = -nFE^\theta_{\text{cell}}\) & 连接热力学与电化学 \\
\end{longtable}
}

\subsection{十、三级练习题}\label{ux5341ux4e09ux7ea7ux7ec3ux4e60ux9898}

\textbf{1.} 判断正误并说明理由：
(a) 放热反应都是自发反应。
(b) 稳定单质的标准生成焓和标准熵均为零。
(c) 状态函数的改变值只与始态和终态有关，与途径无关。
(d) 体系吸热则 Q \textgreater{} 0，所以 ΔH \textgreater{} 0。

\textbf{2.} 写出下列热化学方程式：
(a) 1 mol H\textsubscript{2}(g) 与 ½ mol O\textsubscript{2}(g) 反应生成 H\textsubscript{2}O(g)，放热 241.8 kJ。
(b) 1 mol C(s,石墨) 完全燃烧生成 CO\textsubscript{2}(g)，放热 393.5 kJ。

\textbf{3.} 判断下列反应的 ΔS 符号(正或负)：
(a) NH\textsubscript{4}Cl(s) \箭 NH\textsubscript{3}(g) + HCl(g)
(b) N\textsubscript{2}(g) + 3H\textsubscript{2}(g) \箭 2NH\textsubscript{3}(g)
(c) H\textsubscript{2}O(l) \箭 H\textsubscript{2}O(g)
(d) Ag\textsuperscript{+}(aq) + Cl\textsuperscript{-}(aq) \箭 AgCl(s)

\textbf{4.} 已知 298 K 时 \(\Delta_f H^\theta\) 数据：HCl(g) = −92.3 kJ/mol，H\textsubscript{2}(g) = 0，Cl\textsubscript{2}(g) = 0。计算反应 H\textsubscript{2}(g) + Cl\textsubscript{2}(g) = 2HCl(g) 的 \(\Delta_r H^\theta\)。

\textbf{5.} 某反应 ΔH = −100 kJ/mol，ΔS = −200 J/(mol·K)。判断 298 K 时反应是否自发。若不自发，估算自发所需温度条件。

\textbf{6.} 解释为什么 NaCl 溶于水是自发过程(ΔS \textgreater{} 0，ΔH \textgreater{} 0)，用 ΔG = ΔH − TΔS 说明。

\textbf{7.} 简述 Born-Haber 循环的基本思路，以及它在计算什么物理量时特别有用。

\textbf{8.} 已知某反应的 \(\Delta_r G^\theta\) = −30.0 kJ/mol(298 K)，计算 \(K^\theta\)，并说明反应正向进行的程度。

\textbf{9.} 焓变计算：已知 \(\Delta_f H^\theta\)(H\textsubscript{2}O, l) = −285.8 kJ/mol，\(\Delta_f H^\theta\)(H\textsubscript{2}O, g) = −241.8 kJ/mol。计算 H\textsubscript{2}O(l) \箭 H\textsubscript{2}O(g) 的蒸发焓，并说明其符号的物理意义。

\textbf{10.} 比较以下两组物质的标准摩尔熵大小，并解释原因：
(a) NaCl(s) vs NaCl(aq)
(b) C(s,金刚石) vs C(s,石墨)

\begin{center}\rule{0.5\linewidth}{0.5pt}\end{center}

\textbf{11.} \textbf{(盖斯定律综合)} 已知 298 K 时：
- C(s,石墨) + O\textsubscript{2}(g) = CO\textsubscript{2}(g)，ΔH\textsubscript{1} = −393.5 kJ/mol
- H\textsubscript{2}(g) + ½O\textsubscript{2}(g) = H\textsubscript{2}O(l)，ΔH\textsubscript{2} = −285.8 kJ/mol
- C\textsubscript{2}H\textsubscript{5}OH(l) + 3O\textsubscript{2}(g) = 2CO\textsubscript{2}(g) + 3H\textsubscript{2}O(l)，ΔH\textsubscript{3} = −1367 kJ/mol

求 C\textsubscript{2}H\textsubscript{5}OH(l) 的标准生成焓 \(\Delta_f H^\theta\)。

\textbf{12.} \textbf{(ΔH 与 ΔU 的换算)} 反应 N\textsubscript{2}(g) + 3H\textsubscript{2}(g) = 2NH\textsubscript{3}(g) 在 298 K 时 \(\Delta_r H^\theta\) = −92.2 kJ/mol。计算 \(\Delta_r U^\theta\)，并解释两者差异的物理原因。

\textbf{13.} \textbf{(\(\Delta G^\theta\) 与 \(K^\theta\) 的深度计算)} 已知 298 K 时：
- ½N\textsubscript{2}(g) + ½O\textsubscript{2}(g) = NO(g)，\(\Delta_r G^\theta\) = +86.6 kJ/mol
- 2NO(g) + O\textsubscript{2}(g) = 2NO\textsubscript{2}(g)，\(\Delta_r G^\theta\) = −71.2 kJ/mol

求反应 N\textsubscript{2}(g) + 2O\textsubscript{2}(g) = 2NO\textsubscript{2}(g) 的 \(\Delta_r G^\theta\) 和 \(K^\theta\)。

\textbf{14.} \textbf{(键焓法的局限性)} 用键焓法计算 NH\textsubscript{3}(g) \箭 NH\textsubscript{2}·(g) + H·(g) 所需能量(即 N−H 键的键焓)，并与文献值 389 kJ/mol 比较。已知 NH\textsubscript{3} 中三个 N−H 键的键焓不完全相同，第一、二、三步解离能分别为 435、397、339 kJ/mol。讨论''平均键焓''与''逐级解离能''的区别。

\textbf{15.} \textbf{(Born-Haber 循环)} 利用以下数据计算 KCl 的晶格能：

{\def\LTcaptype{none} % do not increment counter
\begin{longtable}[]{@{}lc@{}}
\toprule\noalign{}
过程 & ΔH (kJ/mol) \\
\midrule\noalign{}
\endhead
\bottomrule\noalign{}
\endlastfoot
K(s) \箭 K(g) & +89 \\
K(g) \箭 K\textsuperscript{+}(g) + e\textsuperscript{-} & +419 \\
½Cl\textsubscript{2}(g) \箭 Cl(g) & +122 \\
Cl(g) + e\textsuperscript{-} \箭 Cl\textsuperscript{-}(g) & −349 \\
K(s) + ½Cl\textsubscript{2}(g) \箭 KCl(s) & −437 \\
\end{longtable}
}

\textbf{16.} \textbf{(转变温度的实际意义)} 已知：
- CaCO\textsubscript{3}(s) = CaO(s) + CO\textsubscript{2}(g)，\(\Delta H^\theta\) = +178 kJ/mol，\(\Delta S^\theta\) = +161 J/(mol·K)
- MgCO\textsubscript{3}(s) = MgO(s) + CO\textsubscript{2}(g)，\(\Delta H^\theta\) = +101 kJ/mol，\(\Delta S^\theta\) = +175 J/(mol·K)

分别估算两者的分解温度，并解释为什么 MgCO\textsubscript{3} 比 CaCO\textsubscript{3} 更容易分解。

\textbf{17.} \textbf{(ΔG 的非标准状态计算)} 在 298 K，某反应的 \(\Delta_r G^\theta\) = +20 kJ/mol。若反应物浓度均为 0.01 mol/L，产物浓度均为 1.0 mol/L，判断反应进行的方向。

\textbf{18.} \textbf{(量热实验设计)} 某同学想用保温杯式量热计测定锌与盐酸反应的焓变。已知：100 mL 1.0 mol/L HCl 加入过量锌粉，温度从 25.0°C 升至 38.2°C。溶液密度 \约等于 1.0 g/mL，比热容 \约等于 4.18 J/(g·K)，量热计常数 C = 50 J/K。
(a) 计算反应放出的热量 Q。
(b) 写出热化学方程式并计算 \(\Delta_r H^\theta\)(以 HCl 为基准)。
(c) 与文献值 \(\Delta_f H^\theta\)(ZnCl\textsubscript{2}, aq) = −485 kJ/mol 比较，讨论可能的误差来源。

\textbf{19.} \textbf{(多步组合)} 利用以下数据，用 Born-Haber 循环计算 NaBr 的晶格能，并与 KCl(第15题结果)比较，讨论离子半径对晶格能的影响。

{\def\LTcaptype{none} % do not increment counter
\begin{longtable}[]{@{}lc@{}}
\toprule\noalign{}
过程 & ΔH (kJ/mol) \\
\midrule\noalign{}
\endhead
\bottomrule\noalign{}
\endlastfoot
Na(s) \箭 Na(g) & +107 \\
Na(g) \箭 Na\textsuperscript{+}(g) + e\textsuperscript{-} & +496 \\
½Br\textsubscript{2}(l) \箭 Br(g) & +112 \\
Br(g) + e\textsuperscript{-} \箭 Br\textsuperscript{-}(g) & −325 \\
Na(s) + ½Br\textsubscript{2}(l) \箭 NaBr(s) & −361 \\
\end{longtable}
}

\begin{center}\rule{0.5\linewidth}{0.5pt}\end{center}

\textbf{20.} \textbf{(第26届全国初赛改编)} 已知 298 K 时：

{\def\LTcaptype{none} % do not increment counter
\begin{longtable}[]{@{}lcc@{}}
\toprule\noalign{}
物质 & \(\Delta_f H^\theta\) (kJ/mol) & \(S^\theta\) (J·mol\textsuperscript{-1}·K\textsuperscript{-1}) \\
\midrule\noalign{}
\endhead
\bottomrule\noalign{}
\endlastfoot
Fe\textsubscript{2}O\textsubscript{3}(s) & −824.2 & 87.4 \\
Al\textsubscript{2}O\textsubscript{3}(s) & −1675.7 & 50.9 \\
Fe(s) & 0 & 27.3 \\
Al(s) & 0 & 28.3 \\
\end{longtable}
}

\begin{enumerate}
\def\labelenumi{(\alph{enumi})}
\tightlist
\item
  计算反应 Fe\textsubscript{2}O\textsubscript{3}(s) + 2Al(s) = Al\textsubscript{2}O\textsubscript{3}(s) + 2Fe(s) 的 \(\Delta_r H^\theta\)、\(\Delta_r S^\theta\) 和 298 K 时的 \(\Delta_r G^\theta\)。
\item
  这就是铝热反应。解释为什么铝热反应一旦引发就极其剧烈(从热力学角度)。
\item
  从 \(\Delta_r G^\theta\) 估算该反应在标准条件下的平衡常数 \(K^\theta\)，说明反应正向进行的程度。
\end{enumerate}

\textbf{21.} \textbf{(第28届全国初赛改编)} 石墨和金刚石在 298 K 时的热力学数据：

{\def\LTcaptype{none} % do not increment counter
\begin{longtable}[]{@{}lcc@{}}
\toprule\noalign{}
物质 & \(\Delta_f H^\theta\) (kJ/mol) & \(S^\theta\) (J·mol\textsuperscript{-1}·K\textsuperscript{-1}) \\
\midrule\noalign{}
\endhead
\bottomrule\noalign{}
\endlastfoot
C(s,石墨) & 0 & 5.74 \\
C(s,金刚石) & +1.89 & 2.38 \\
\end{longtable}
}

\begin{enumerate}
\def\labelenumi{(\alph{enumi})}
\tightlist
\item
  计算 C(石墨) \箭 C(金刚石) 的 \(\Delta_r G^\theta\)。
\item
  从热力学角度解释为什么自然界中金刚石稀少而石墨常见。
\item
  已知金刚石密度 = 3.51 g/cm\textsuperscript{3}，石墨密度 = 2.26 g/cm\textsuperscript{3}。施加高压有利于向哪个方向转化？利用 \(\left(\frac{\partial G}{\partial p}\right)_T = V\) 定性分析。
\end{enumerate}

\textbf{22.} \textbf{(全国初赛)} 某反应在 298 K 时 \(\Delta_r H^\theta\) = −120 kJ/mol，\(\Delta_r S^\theta\) = −150 J/(mol·K)。
(a) 判断 298 K 时反应是否自发。
(b) 估算反应自发进行的温度上限。
(c) 若希望该反应在 500 K 时也能自发进行，可以通过什么方式实现？(列举两种方案)

\textbf{23.} \textbf{(综合计算)} 已知 298 K 时：

{\def\LTcaptype{none} % do not increment counter
\begin{longtable}[]{@{}lcc@{}}
\toprule\noalign{}
物质 & \(\Delta_f H^\theta\) (kJ/mol) & \(S^\theta\) (J·mol\textsuperscript{-1}·K\textsuperscript{-1}) \\
\midrule\noalign{}
\endhead
\bottomrule\noalign{}
\endlastfoot
H\textsubscript{2}O\textsubscript{2}(l) & −187.8 & 109.6 \\
H\textsubscript{2}O(l) & −285.8 & 69.9 \\
O\textsubscript{2}(g) & 0 & 205.1 \\
\end{longtable}
}

\begin{enumerate}
\def\labelenumi{(\alph{enumi})}
\tightlist
\item
  计算 2H\textsubscript{2}O\textsubscript{2}(l) = 2H\textsubscript{2}O(l) + O\textsubscript{2}(g) 的 \(\Delta_r H^\theta\)、\(\Delta_r S^\theta\)、\(\Delta_r G^\theta\)。
\item
  计算 298 K 时的 \(K^\theta\)。
\item
  解释为什么 H\textsubscript{2}O\textsubscript{2} 溶液在室温下会缓慢分解，但纯 H\textsubscript{2}O\textsubscript{2} 可以稳定存在。
\item
  从动力学角度补充说明：虽然热力学预测反应自发，但实际分解速率取决于什么？
\end{enumerate}

\textbf{24.} \textbf{(电化学桥接题)} 298 K 时，Zn \textbar{} Zn\textsuperscript{2+}(1 mol/L) \textbar\textbar{} Cu\textsuperscript{2+}(1 mol/L) \textbar{} Cu 电池的标准电动势 \(E^\theta\) = 1.10 V。
(a) 用 \(\Delta_r G^\theta = -nFE^\theta\) 计算该电池反应的 \(\Delta_r G^\theta\)。
(b) 用 \(\Delta_r G^\theta = -RT \ln K^\theta\) 计算反应 Zn(s) + Cu\textsuperscript{2+}(aq) = Zn\textsuperscript{2+}(aq) + Cu(s) 的 \(K^\theta\)。
(c) 解释为什么该电池反应的 \(K^\theta\) 如此之大(\textgreater{} 10\textsuperscript{37})。

\textbf{25.} \textbf{(综合推断)} 某化合物 X 在加热时发生分解反应，产生两种气体。已知：
- 分解温度约为 450 K(此时 \(K^\theta\) \约等于 1)
- 分解反应是吸热反应
- 产物气体分子数多于反应物

\begin{enumerate}
\def\labelenumi{(\alph{enumi})}
\tightlist
\item
  推测 X 可能是什么类型的化合物(给出 2−3 种可能)。
\item
  从 ΔG = ΔH − TΔS 分析，为什么分解温度是 450 K 而不是更低？
\item
  若对体系施加高压，分解温度会如何变化？用勒夏特列原理和热力学分析。
\end{enumerate}

\subsection{十一、练习题答案}\label{ux5341ux4e00ux7ec3ux4e60ux9898ux7b54ux6848}

\textbf{1.}
(a) 错误。放热(ΔH \textless{} 0)\不等于 自发。自发性由 ΔG = ΔH − TΔS 决定。若 ΔS \textless{} 0 且温度足够高，ΔG \textgreater{} 0 可使放热反应非自发。
(b) 错误。稳定单质的 \(\Delta_f H^\theta\) = 0，但 \(S^\theta\) \不等于 0。只有在 0 K 时完美晶体的熵才为零。
(c) 正确。这正是状态函数的定义特征。
(d) 错误。吸热 Q \textgreater{} 0 时，ΔH = Qp \textgreater{} 0 仅在恒压条件下成立。且 ΔH 与 Q 的取号规则不同。

\textbf{2.}
(a) H\textsubscript{2}(g) + ½O\textsubscript{2}(g) = H\textsubscript{2}O(g)，\(\Delta_r H^\theta\) = −241.8 kJ/mol
(b) C(s,石墨) + O\textsubscript{2}(g) = CO\textsubscript{2}(g)，\(\Delta_r H^\theta\) = −393.5 kJ/mol

\textbf{3.}
(a) ΔS \textgreater{} 0(固体\箭气体，混乱度大增)
(b) ΔS \textless{} 0(4 mol 气体\箭2 mol 气体，分子数减少)
(c) ΔS \textgreater{} 0(液体\箭气体)
(d) ΔS \textless{} 0(离子\箭固体沉淀，有序度增加)

\textbf{4.} \(\Delta_r H^\theta\) = 2 × (−92.3) − {[}0 + 0{]} = −184.6 kJ/mol

\textbf{5.} ΔG = −100 − 298 × (−0.200) = −100 + 59.6 = −40.4 kJ/mol \textless{} 0，反应在 298 K 时自发。若在更高温度(如 T \textgreater{} 500 K)，−TΔS 项变大正值，ΔG 可能变为正值。

\textbf{6.} NaCl 溶解：ΔH \textgreater{} 0(吸热)，ΔS \textgreater{} 0(离子分散到溶液中，混乱度增加)。ΔG = ΔH − TΔS，室温下 TΔS 项足够大，使 ΔG \textless{} 0，所以自发。

\textbf{7.} Born-Haber 循环将离子晶体的生成过程拆解为升华\箭电离\箭解离\箭电子亲和\箭晶格形成等步骤。由盖斯定律，各步骤焓变之和等于生成焓。特别适用于间接求晶格能。

\textbf{8.} \(K^\theta\) = e\^{}(30000/(8.314×298)) = e\^{}12.1 \约等于 1.8 × 10\textsuperscript{5}。\(K^\theta\) 远大于 1，反应正向进行很彻底。

\textbf{9.} \(\Delta_r H^\theta\) = −241.8 − (−285.8) = +44.0 kJ/mol。正值表示水的蒸发是吸热过程。

\textbf{10.}
(a) \(\mathrm{NaCl(aq)}\) 的 \(S^\theta\) \textgreater{} \(\mathrm{NaCl(s)}\) 的 \(S^\theta\)：溶解后离子分散到溶液中，混乱度增大。
(b) 石墨的 \(S^\theta\) \textgreater{} 金刚石的 \(S^\theta\)：石墨为层状结构，层间有自由电子运动，微观状态数更多。

\begin{center}\rule{0.5\linewidth}{0.5pt}\end{center}

\textbf{11.} 目标反应 = 2×(1) + 3×(2) − (3)
\(\Delta_f H^\theta\) = 2×(−393.5) + 3×(−285.8) − (−1367) = −277.4 kJ/mol

\textbf{12.} Δn = 2 − 4 = −2，\(\Delta_r U^\theta\) = \(\Delta_r H^\theta\) − ΔnRT = −92.2 − (−2) × 8.314 × 10\textsuperscript{-3} × 298 = −87.2 kJ/mol。物理原因：反应使气体分子数减少，环境对体系做功。

\textbf{13.} \(\Delta_r G^\theta\) = 2×(+86.6) + (−71.2) = +102.0 kJ/mol；\(K^\theta\) = e\^{}(-102000/2478) \约等于 1.2 × 10\textsuperscript{-18}。\(K^\theta\) 极小，N\textsubscript{2} 与 O\textsubscript{2} 直接生成 NO\textsubscript{2} 几乎不可能。

\textbf{14.} 平均键焓 = (435+397+339)/3 = 390 kJ/mol，与文献值 389 kJ/mol 接近。但逐级解离能差异显著：第一个 N−H 键最难断裂(435 kJ)，最后一个最易(339 kJ)。``平均键焓''掩盖了这种逐级差异。

\textbf{15.} U = −437 − 89 − 419 − 122 + 349 = −718 kJ/mol(绝对值 718 kJ/mol)

\textbf{16.} CaCO\textsubscript{3}：T\_转 \约等于 1106 K；MgCO\textsubscript{3}：T\_转 \约等于 577 K。MgCO\textsubscript{3} 更易分解，因为 MgO 晶格能更大(Mg\textsuperscript{2+} 半径更小)。

\textbf{17.} Q = (1.0)\textsuperscript{2}/(0.01)\textsuperscript{2} = 10\textsuperscript{4}，ΔG = +20000 + 8.314×298×ln(10\textsuperscript{4}) = +42.8 kJ/mol \textgreater{} 0，反应逆向自发。

\textbf{18.}
(a) Q = (4.18×100×1.0+50)×(38.2−25.0) = 468×13.2 = 6178 J \约等于 6.18 kJ
(b) n(HCl) = 0.100 mol，\(\Delta_r H^\theta\) = −6.18/0.100 = −61.8 kJ/mol
(c) 文献值 −485 kJ/mol，偏差较大。可能误差：H\textsubscript{2} 气体带走热量、保温不理想、温度读数误差。

\textbf{19.} U(NaBr) = −361 − 107 − 496 − 112 + 325 = −751 kJ/mol。NaBr 晶格能更大(751 \textgreater{} 718)，因为 Na\textsuperscript{+} 和 Br\textsuperscript{-} 的半径之和比 K\textsuperscript{+} 和 Cl\textsuperscript{-} 更小。

\begin{center}\rule{0.5\linewidth}{0.5pt}\end{center}

\textbf{20.}
(a) \(\Delta_r H^\theta\) = −851.5 kJ/mol，\(\Delta_r S^\theta\) = −31.9 J/(mol·K)，\(\Delta_r G^\theta\) = −842.0 kJ/mol
(b) \(\Delta_r G^\theta\) 极负，产物热力学稳定性远高于反应物。一旦引发，大量焓变以热能释放，温度急剧升高---典型的''热爆炸''式自催化过程。
(c) \(K^\theta\) \约等于 10\^{}147，反应正向极其彻底。

\textbf{21.}
(a) \(\Delta_r G^\theta\) = +2.89 kJ/mol \textgreater{} 0
(b) \(\Delta_r G^\theta\) \textgreater{} 0，石墨\箭金刚石非自发。石墨是热力学稳定相。
(c) 高压有利于向密度更大的相转化((∂G/∂p)\_T = V)。金刚石密度 \textgreater{} 石墨，施压有利于石墨\箭金刚石。

\textbf{22.}
(a) \(\Delta G^\theta\) = −120 + 44.7 = −75.3 kJ/mol \textless{} 0，自发。
(b) T\_转 = 800 K。
(c) 降低温度；改变浓度/分压使 Q \textless{} K；耦合一个 ΔG 更负的反应。

\textbf{23.}
(a) \(\Delta_r H^\theta\) = −196.0 kJ/mol，\(\Delta_r S^\theta\) = +125.7 J/(mol·K)，\(\Delta_r G^\theta\) = −233.5 kJ/mol
(b) \(K^\theta\) \约等于 10\^{}41
(c) 热力学上 \(K^\theta\) 极大。H\textsubscript{2}O\textsubscript{2} 溶液中微量杂质催化分解；纯 H\textsubscript{2}O\textsubscript{2} 缺乏催化活性位点。
(d) 分解反应活化能较高(O−O 键断裂需 \textasciitilde210 kJ/mol)，室温下分子动能不足以大量跨越能垒。

\textbf{24.}
(a) \(\Delta_r G^\theta\) = −2 × 96485 × 1.10 = −212.3 kJ/mol
(b) \(K^\theta\) = e\^{}(212300/2478) \约等于 10\^{}37
(c) Cu\textsuperscript{2+}/Cu 和 Zn\textsuperscript{2+}/Zn 之间的电极电势差很大(1.10 V)，电子转移的热力学驱动力极大。

\textbf{25.}
(a) 可能：CaCO\textsubscript{3}、NH\textsubscript{4}HCO\textsubscript{3}、CuSO\textsubscript{4}·5H\textsubscript{2}O 等。
(b) 450 K 时 \(K^\theta\) \约等于 1 即 \(\Delta G^\theta\) \约等于 0。T \textless{} 450 K 时 \(T\Delta S^\theta < \Delta H^\theta\)，\(\Delta G^\theta\) \textgreater{} 0；T \textgreater{} 450 K 时 \(\Delta G^\theta\) \textless{} 0。
(c) 加压使分解温度升高。分解产生气体(体积增大)，加压不利于气体产物生成。

\end{document}
